\documentclass{article}
\usepackage{amsmath,amssymb,amsthm}
\newtheorem{theorem}{Theorem}
\title{Weak-coupling resummation of one complete conditional Wilson cube}\author{}\date{}
\begin{document}\maketitle
This is a real, finite-block weak-coupling estimate for the actual Wilson conditional law. Unlike the six-face primitive in the exact cube scalar obstruction calculation, it retains every exterior plaquette incident to the integrated links. It retains the full scalar normalization. It is not an all-history R transition or a uniform estimate over arbitrary baths.

\section{Exact conditional record and its flat Hessian}
Embed one elementary three-dimensional cube in the four-dimensional lattice, with its fourth coordinate fixed. Take its twelve positively oriented edges as the free set $E_i$, and condition on every other link. The torus is large enough that this record and its one-plaquette collar do not wrap or identify vertices. Denote the fixed bath by $b$.

There are six plaquettes entirely in the free cube. Each free edge belongs to exactly four additional plaquettes, each containing no other free edge. Thus the complete action depending on $E_i$ has exactly $6+48=54$ plaquettes:
\[
 S_b(U)=\sum_{P\in\partial\text{cube}}s(U_P)
 +\sum_{e\in E_i}\sum_{r=1}^{4}s(U_eV_{e,r}(b)).
 \tag{1}
\]
Inverses and orientation changes are absorbed using $s(W^{-1})=s(W)$ and cyclicity of trace; each $V_{e,r}$ is the original ordered three-link bath staple, with its original endpoints. The denominator and conditional law are exactly
\[
 D_\beta(b)=\int_{G^{12}}e^{-\beta S_b(U)}dU,
 \qquad K_{\beta,b}(dU)=D_\beta(b)^{-1}e^{-\beta S_b(U)}dU.
 \tag{2}
\]
The remaining action of the full lattice is independent of these twelve links and factors outside (2).

At the literal flat bath $b=I$, all $V_{e,r}=I$, and
\[
 S_I(U)=\sum_{P\in\partial\text{cube}}s(U_P)+4\sum_{e\in E_i}s(U_e).
 \tag{3}
\]
This has a unique global minimum $U=I$. Use an orthonormal basis of $\mathfrak{su}(N)$ for the Frobenius metric, so $s(e^A)=|A|^2/(2N)+O(|A|^4)$. Let $d$ be the oriented cube face-edge incidence matrix. The full Hessian in all $n_i=12d_G$ free coordinates is
\[
 H_I=N^{-1}(4I_{12}+d^Td)\otimes I_{d_G},\qquad d_G=N^2-1.
 \tag{4}
\]
The six-by-six matrix $dd^T$ has eigenvalues $0,4,4,4,6,6$. Hence $H_I$ has eigenvalues $4/N$ (multiplicity $7d_G$), $8/N$ (multiplicity $3d_G$), and $10/N$ (multiplicity $2d_G$), and
\[
 \det H_I=N^{-12d_G}(4^7 8^3 10^2)^{d_G}.
 \tag{5}
\]
The seven would-be tree/gauge directions of the isolated cube are not divided out: the fixed exterior staples pin them and give the eigenvalue $4/N$. There are twelve original Haar integrations in (2), not five. The five-integration count in \texttt{exact\_cube\_scalar\_obstruction.tex} applies only to the isolated closed-surface factor, which has a different boundary condition.

\section{Uniform nearby-bath Laplace estimate with scalar retained}
Let $J_0$ denote the density of normalized Haar measure at zero in these fixed orthonormal exponential coordinates. It is kept explicitly.
\begin{theorem}
There is a fixed open neighborhood $\mathcal B$ of the flat complete bath record and constants independent of $\beta$ such that, for every $b$ in a sufficiently small closed sub-neighborhood, the action (1) has a unique global minimum $U_*(b)$, positive Hessian $H_b\geq cI_{n_i}$, and
\[
 D_\beta(b)=e^{-\beta S_b(U_*)}
 J_0^{12}(2\pi/\beta)^{n_i/2}(\det H_b)^{-1/2}
 [1+\varepsilon_\beta(b)],\qquad
 |\varepsilon_\beta(b)|\leq C\beta^{-1/2}.
 \tag{6}
\]
In particular
\[
 \log D_\beta(b)=-\beta S_b(U_*)-\frac{n_i}{2}\log\beta
 +12\log J_0+\frac{n_i}{2}\log(2\pi)
 -\frac12\log\det H_b+O(\beta^{-1/2}).
 \tag{7}
\]
In the coordinates $U_e=U_{*,e}(b)e^{x_e/\sqrt\beta}$, the normalized law converges to the full Gaussian with covariance $H_b^{-1}$. For each fixed integer $p\geq0$, there is $\theta>0$ such that its density $k_{\beta,b}$ on the fixed logarithm chart obeys
\[
 \int_{\mathbb R^{n_i}}(1+|x|)^pe^{\theta|x|^2}
 |k_{\beta,b}(x)-g_{H_b}(x)|dx\leq C_p\beta^{-1/2},
 \tag{8}
\]
where $k$ is extended by zero outside the rescaled chart. The omitted probability is at most $Ce^{-c\beta}$. All assertions are uniform over that nearby-bath set.
\end{theorem}
\begin{proof}
At $b=I$ the compact action (3) has a unique zero and a strictly positive Hessian (4). In a small convex coordinate ball its Hessian remains uniformly positive. The implicit function theorem supplies a unique critical point there for nearby $b$. On the compact complement of a smaller ball, the flat action has a positive gap. Uniform continuity preserves a positive gap relative to the nearby critical value after the bath neighborhood is reduced. This proves global uniqueness, uniform Hessian bounds, and the exponential exit estimate.

Translate coordinates by left multiplication with $U_*(b)$. Haar invariance keeps the local density at zero equal to $J_0^{12}$. Taylor's theorem, uniformly on the compact bath set, gives
\[
 \beta[S_b(U_*e^{x/\sqrt\beta})-S_b(U_*)]
 =\tfrac12 x^TH_bx+R_{\beta,b}(x),\quad
 |R_{\beta,b}(x)|\leq C\beta^{-1/2}|x|^3.
\]
On the original sufficiently small chart both the exact action difference and its quadratic approximation are bounded below by $c|x|^2/\beta$. The mean-value formula for two exponentials therefore bounds their difference by
$C\beta^{-1/2}|x|^3e^{-c'|x|^2}$, reducing the chart if necessary. The smooth Haar Jacobian differs from $J_0^{12}$ by at most $C\beta^{-1/2}(1+|x|)$ there; the sharper quadratic estimate is unnecessary. Integration proves (6). Taking a logarithm proves (7). Division by the same genuine denominator in (6), followed by a Gaussian weighted integral with $\theta<c'$, proves (8). Outside the chart the global gap absorbs the fixed power $\beta^{n_i/2}$. No uncounted gauge volume or independent scalar is removed.
\end{proof}

At the flat bath, (7) and (5) are an explicit scalar resummation with $n_i/2=6d_G$. The law is weak-coupling small relative to its full pinned Gaussian, not relative to independent Haar boundary variables. Gauge-equivalent copies of the bath give identical conclusions by transforming all fixed staples and free links together. This does not assert that a chosen original R gauge or its selector lands in this bath neighborhood.

\section{What survives in the complete lattice after integration}
Integrating this block replaces its factor by $D_\beta(b)$. Thus the effective bath action contains the exact local increment $-\log D_\beta(b)$, including the minimum term, $6d_G\log\beta$, and the Hessian determinant. The normalized error in (8) does not bound this scalar by $O(\beta^{-1/2})$.

There is also a precise reason the pinned mass in (4) cannot become a global mass. At the flat full-lattice quadratic level split all links into the cube coordinates $i$ and bath coordinates $b$, writing the curl Hessian as $Q$. The induced bath quadratic form is the Schur complement
\[
 Q_{\rm eff}=Q_{bb}-Q_{bi}Q_{ii}^{-1}Q_{ib}.
 \tag{9}
\]
If $Q(v_i,v_b)=0$, its first block row gives $v_i=-Q_{ii}^{-1}Q_{ib}v_b$, and its second gives $Q_{\rm eff}v_b=0$. Thus the exact global infinitesimal gauge and torus harmonic modes are retained. Conditioning pins a finite block; integrating it does not eliminate the collective modes or make the remaining bath interaction small.

For genuinely copied complete cube-and-staple records under a torus covering, (1)--(2) show pointwise equality of $D_\beta$, $U_*$ (in correspondingly transported charts), and all terms of (7). A family of disjoint free blocks with no plaquette coupling different blocks factors conditionally, so this pointwise scalar matching has exactly the covering multiplicity. This is a useful local full-record comparison. It is not a comparison of the distributions of unmatched baths, an identification of arbitrary R histories, or control of wrapped records.

The remaining S94 estimate is now explicit: control or match the induced bath action and its harmonic/boundary modes, with the scalar (7) retained. Uniform admission of these blocks under the manuscript's original selectors and the all-scale sum have not been proved.
\end{document}
